\subsection*{J5: the height formula fails for a split quadric fourfold}
\textbf{Status: Refuted.} Root independently audited the complete proof and both supplied transcriptions of 44:4.
The failure below concerns the literal statement of Theorem5, 44:4.
Its insertion reads ``section hyperplane générique d'un pinceau dont
[les] fibres sont géom. intègres'', and it concludes
$\langle\zeta,\zeta'\rangle=(-1)^mZ\mathbin{.}Z'$.
The example satisfies this statement,
including the inserted requirement that all pencil fibres be
geometrically integral. It does not depend on identifying the
manuscript's equivalence relation or constructing any nonzero height:
both restricted cycles are literally zero. The original middle
case $n=2m=4$, $i=j=m=2$ also makes its displayed sign unambiguous.

Let $k=\overline{\mathbb F}_3$, and let
\[
 X=V(x_0x_3+x_1x_4+x_2x_5)\subset\mathbb P^5_k.
\]
The six partial derivatives are $x_3,x_4,x_5,x_0,x_1,x_2$,
so $X$ is a smooth projective fourfold. Consider the pencil
\[
 H_t=V(x_2-tx_5),\quad t\in\mathbb A^1,
 \qquad H_\infty=V(x_5).
\]
Its base locus in $X$ is the smooth quadric surface
$B=V(x_2,x_5,x_0x_3+x_1x_4)\subset\mathbb P^5$.
Blowing up $B$ resolves the pencil and gives a smooth projective
total space over $\mathbb P^1$; the fibres remain $X\cap H_t$,
because $B$ is a Cartier divisor in each member. For finite $t$,
the fibre has equation
\[
                 x_0x_3+x_1x_4+t x_5^2=0
                 \quad\text{in }\mathbb P^4.
\]
It is smooth for $t\ne0$ and is a rank-four quadric cone for $t=0$.
The fibre at infinity is another rank-four quadric cone, with
equation $x_0x_3+x_1x_4=0$ and extra coordinate $x_2$.
All fibres are geometrically integral: a quadratic form which is
a product of two linear forms has rank at most two, whereas these
have rank at least four, over every algebraic extension.
In particular this is a Lefschetz pencil with smooth base and two
ordinary quadratic singular fibres, rather than a pencil whose
total family is degenerate.

Take the four planes
\[
 \alpha=V(x_3,x_4,x_5)\subset X,
 \qquad
 \beta=V(x_2,x_3,x_4)\subset X,
\]
\[
 \alpha'=V(x_0,x_1,x_5)\subset X,
 \qquad
 \beta'=V(x_0,x_1,x_2)\subset X.
\]
Put $Z=[\alpha]-[\beta]$, $Z'=[\alpha']-[\beta']$.
If $K=k(t)$ and $Y=X\cap H_t$ is the generic
member, then the proper scheme-theoretic intersections are
\[
 \alpha\cap Y=\beta\cap Y
   =L:=V(x_2,x_3,x_4,x_5)\simeq\mathbb P^1_K,
\]
and
\[
 \alpha'\cap Y=\beta'\cap Y
   =L':=V(x_0,x_1,x_2,x_5)\simeq\mathbb P^1_K,
\]
all with multiplicity one. Indeed $H_t$ restricts on $\alpha$
to $x_2=0$, and on $\beta$ to $-t x_5=0$, where $t$ is a unit
in $K$; the same statements hold for $\alpha',\beta'$.
Therefore
\[
                  \zeta=Z\cdot Y=0,
                  \qquad \zeta'=Z'\cdot Y=0             \tag{J5.1}
\]
already in the group of cycles, not only modulo an equivalence.
Any Neron--Tate bilinear pairing on their associated points is zero.

Here is a direct calculation using only proper transverse
intersections. The crossed pairs $\alpha\cap\beta'$ and
$\beta\cap\alpha'$ are empty. The other two intersections are
\[
 \alpha\cap\alpha'=\{[e_2]\},\qquad
 \beta\cap\beta'=\{[e_5]\}.
\]
At $[e_2]$, the chart $x_2=1$ writes the smooth fourfold with local
coordinates $x_0,x_1,x_3,x_4$ and
$x_5=-x_0x_3-x_1x_4$. The ideals of the two planes are
$(x_3,x_4)$ and $(x_0,x_1)$. They therefore intersect transversely
with local multiplicity one. At $[e_5]$, the chart $x_5=1$
gives exactly the same coordinate calculation, with
$x_2=-x_0x_3-x_1x_4$. Consequently
\[
                  Z\cdot Z'=1+1-0-0=2.                 \tag{J5.2}
\]
The integer two in (J5.2) is the degree of a zero-cycle; it is not
to be reduced modulo the characteristic.

Equations (J5.1)--(J5.2) contradict exactly the proposed formula
\[
             \langle\zeta,\zeta'\rangle
                       =(-1)^m Z\cdot Z',
                  \qquad m=2,
\]
by giving $0=2$. Changing one overall sign cannot fix it.
Both cycles have algebraically trivial generic restrictions, and
every member of the pencil is geometrically integral, as required
by the revised ink. Thus the literal higher-codimension Theorem5
is false even in this smallest balanced case beyond surfaces.

\paragraph{What the integral-fibre condition misses.}
The obstruction is exact: generic restriction does not determine
intersection numbers of higher-codimension extensions. The class
$Z$ is in the kernel of generic restriction but pairs nontrivially with $Z'$,
so this intersection pairing cannot descend to that quotient.
In this displayed example each plane meets the pencil base in a line,
so those intersections have excess dimension. The literal statement
does not forbid this. An additional condition on extensions or their
incidence with the base would be a stronger assertion requiring a
separate analysis. After blowing up the pencil base, the strict transform of $\alpha$
lies in the fibre at infinity and that of $\beta$ in the fibre at
zero; their common exceptional correction cancels in the difference.
Geometric integrality eliminates independent vertical divisor
components, but it does not eliminate these vertical codimension-two
classes through the singularities of the quadric fibres.

The following precise elementary replacement records the necessary
correction. For any rational vector space $V$ with a bilinear form
$b$ and a restriction map $r:V\to W$, let $K=\ker r$.
The form factors through $r(V)$ if and only if
$b(K,V)=b(V,K)=0$. Necessity follows by pairing a zero image with
every image; sufficiency follows by changing either lift by an
element of $K$. For a symmetric form, its restriction always
descends to
\[
                K^\perp/(K\cap K^\perp).
\]
Apply this to rational cycle classes on a smooth projective model,
their intersection form, and restriction to the generic fibre.
A higher-codimension height comparison needs an admissibility or
vertical-correction condition that supplies this orthogonality;
geometric integrality of the fibres alone does not supply it.
This counterexample does not refute a comparison with such
additional conditions, nor the classical divisor comparison for
surfaces with integral fibres.

\paragraph{Finite check.}
The saved script \texttt{check\_quadric\_height.py} checks the
gradient, the pencil ranks, all four literal restriction ideals,
the two transverse point intersections and the two empty intersections,
and the integer intersection value. An additional normal-bundle and
moving-plane calculation independently checks $Z^2=2$ as well.
The written proof supplies the geometric and height conclusions.
